\documentclass[11pt]{article}
% ===========================================================================
%  Shared preamble for the AP Calculus AB exam-preparation worksheets.
%  Input right after \documentclass. Files beginning with "_" are not compiled
%  on their own. Navy and gold, to match the course's lesson pages.
% ===========================================================================
\usepackage[letterpaper,top=2.3cm,bottom=2.2cm,left=1.7cm,right=1.7cm,headheight=28pt,footskip=30pt]{geometry}
\usepackage{amsmath,amssymb}
\usepackage[T1]{fontenc}
\usepackage[utf8]{inputenc}
\usepackage{lmodern}
\usepackage{xcolor}
\usepackage[most]{tcolorbox}
\usepackage{enumitem}
\usepackage{fancyhdr}
\usepackage{array}
\usepackage{tikz}
\usepackage{pgfplots}
\pgfplotsset{compat=1.18}

\definecolor{navy}{HTML}{161E54}
\definecolor{ink}{HTML}{27357E}
\definecolor{gold}{HTML}{B98A00}
\definecolor{goldbg}{HTML}{FFFDF3}
\definecolor{inkbg}{HTML}{F6F7FD}
\definecolor{rule}{HTML}{C9CFF0}
\definecolor{okgreen}{HTML}{1B7A44}
\definecolor{okbg}{HTML}{F1F8F3}

\pagestyle{fancy}
\fancyhf{}
\fancyhead[L]{\footnotesize Name: \underline{\hspace{4cm}}}
\fancyhead[C]{\textbf{\color{navy}Integration Academy}\\[-2pt]{\scriptsize Calculus AB \textperiodcentered{} AP\textsuperscript{\textregistered} Exam Preparation}}
\fancyfoot[L]{\tiny\color{black!55}\parbox[t]{0.78\textwidth}{AP\textsuperscript{\textregistered} is a registered trademark of the College Board, which was not involved in the production of, and does not endorse, this material. \textcopyright{} Integration Academy.}}
\fancyfoot[R]{\scriptsize\thepage}
\renewcommand{\headrulewidth}{1.2pt}
\renewcommand{\headrule}{\hbox to\headwidth{\color{gold}\leaders\hrule height \headrulewidth\hfill}}
\renewcommand{\footrulewidth}{0pt}
\setlength{\parindent}{0pt}
\setlength{\parskip}{4pt}

% Title block: \wstitle{code}{title}{unit line}
\newcommand{\wstitle}[3]{\begin{center}{\small\bfseries\color{gold}WORKSHEET #1}\\[3pt]{\LARGE\bfseries\color{navy}#2}\\[3pt]{\small\color{black!60}#3}\end{center}\vspace{2pt}}
% Section heading
\newcommand{\sech}[1]{\par\vspace{10pt}{\large\bfseries\color{navy}#1}\par\nobreak\vspace{2pt}{\color{rule}\hrule height0.8pt}\vspace{7pt}}
% A level tag printed beside a group of questions
\newcommand{\level}[1]{\par\vspace{6pt}{\footnotesize\bfseries\color{gold}\MakeUppercase{#1}}\par\nobreak\vspace{2pt}}

\newtcolorbox{ideabox}{colframe=ink,colback=inkbg,coltitle=white,colbacktitle=ink,fonttitle=\bfseries\small,title={Key ideas},arc=2pt,boxrule=0.8pt,left=8pt,right=8pt}
\newtcolorbox{example}[1]{colframe=rule,colback=white,coltitle=navy,colbacktitle=inkbg,fonttitle=\bfseries\small,title={#1},arc=2pt,boxrule=0.8pt,borderline west={3pt}{0pt}{ink},left=10pt,right=8pt}
\newtcolorbox{frbox}[1]{colframe=gold,colback=goldbg,coltitle=white,colbacktitle=gold,fonttitle=\bfseries\small,title={#1},arc=2pt,boxrule=0.8pt,breakable,left=8pt,right=8pt}
\newtcolorbox{solbox}{colframe=okgreen,colback=okbg,coltitle=white,colbacktitle=okgreen,fonttitle=\bfseries\small,title={Worked solutions},arc=2pt,boxrule=0.8pt,breakable,left=8pt,right=8pt}
\newtcolorbox{scorebox}[1]{colframe=navy,colback=white,coltitle=white,colbacktitle=navy,fonttitle=\bfseries\small,title={#1},arc=2pt,boxrule=0.8pt,breakable,left=8pt,right=8pt}

\newcommand{\soln}{\par\smallskip\textbf{\color{ink}Solution.}\ }
\newcommand{\result}[1]{\par\smallskip\textbf{#1}}
% blank working space of a given height
\newcommand{\work}[1]{\par\vspace{#1}}
% a point value in the margin of a scoring guide
\newcommand{\pts}[1]{\hfill{\footnotesize\bfseries\color{gold}[#1]}}

% Coordinate grid for graph questions: \begin{apgraph}{xmin}{xmax}{ymin}{ymax} ... \end{apgraph}
\newenvironment{apgraph}[4]{\begin{center}\begin{tikzpicture}\begin{axis}[
  axis lines=middle, xmin=#1, xmax=#2, ymin=#3, ymax=#4,
  grid=both, grid style={black!12}, xtick distance=1, ytick distance=1,
  ticklabel style={font=\scriptsize}, xlabel={$x$}, ylabel={$y$},
  width=10cm, height=7.2cm, samples=80, clip=false,
  every axis plot/.append style={line width=1.1pt, color=ink}]}{\end{axis}\end{tikzpicture}\end{center}}
\newcommand{\closeddot}[2]{\addplot[only marks, mark=*, mark size=2.2pt, color=ink] coordinates {(#1,#2)};}
\newcommand{\opendot}[2]{\addplot[only marks, mark=*, mark size=2.2pt, color=ink, mark options={fill=white}] coordinates {(#1,#2)};}

% Multiple-choice options: four across, or stacked when the options are long.
\newcommand{\opts}[4]{\par\vspace{3pt}\begin{tabular}{@{}p{0.235\linewidth}p{0.235\linewidth}p{0.235\linewidth}p{0.235\linewidth}@{}}(A)\ #1 & (B)\ #2 & (C)\ #3 & (D)\ #4\end{tabular}\par\vspace{9pt}}
\newcommand{\optsv}[4]{\par\vspace{3pt}\begin{tabular}{@{}l@{\ \ }p{0.9\linewidth}@{}}(A) & #1\\(B) & #2\\(C) & #3\\(D) & #4\end{tabular}\par\vspace{9pt}}

\fancyhead[R]{\footnotesize\color{black!60}Worksheet 6.11}
\begin{document}

\wstitle{6.11}{Choosing an Integration Method}{Unit 6: Integration and Accumulation of Change}

\begin{ideabox}
\begin{itemize}[leftmargin=*,itemsep=2pt,topsep=2pt]
  \item First ask whether \textbf{algebra} is enough: expand a product, or divide a quotient term by term.
  \item Then look for a \textbf{substitution}: an inside function whose derivative is present, up to a constant.
  \item A numerator that is the derivative of the denominator gives a \textbf{logarithm}.
  \item A constant over a quadratic with no real zeros gives an \textbf{arctangent}, after completing the square.
\end{itemize}
\end{ideabox}

\sech{Worked examples}
\begin{example}{Example 1 \textperiodcentered{} A logarithm}
Find $\displaystyle\int\dfrac{2x+1}{x^2+x+3}\,dx$.\soln The numerator is the derivative of the denominator.\result{$\ln\!\left(x^2+x+3\right)+C$.}
\end{example}

\begin{example}{Example 2 \textperiodcentered{} Algebra first}
Find $\displaystyle\int\dfrac{(x+1)^2}{\sqrt{x}}\,dx$.\soln Expand and divide: $x^{3/2}+2x^{1/2}+x^{-1/2}$.\result{$\dfrac25x^{5/2}+\dfrac43x^{3/2}+2\sqrt{x}+C$.}
\end{example}

\begin{example}{Example 3 \textperiodcentered{} An arctangent}
Find $\displaystyle\int\dfrac{1}{x^2+2x+2}\,dx$.\soln $x^2+2x+2=(x+1)^2+1$.\result{$\arctan(x+1)+C$.}
\end{example}

\newpage
\sech{Practice}

\level{Core}
\begin{enumerate}[leftmargin=*,itemsep=2pt]
  \item Find $\displaystyle\int\left(3x^2+1\right)\left(x^3+x\right)^5dx$.\work{2.3cm}
  \item Find $\displaystyle\int\dfrac{x^2+1}{x^2}\,dx$.\work{2.3cm}
  \item Find $\displaystyle\int\dfrac{1}{1+4x^2}\,dx$.\work{2.3cm}
  \item Find $\displaystyle\int\dfrac{e^{x}}{e^{x}+2}\,dx$.\work{2.3cm}
\end{enumerate}

\level{Exam level}
\begin{enumerate}[leftmargin=*,itemsep=2pt,start=5]
  \item Find $\displaystyle\int\dfrac{\sin x}{\cos^2x}\,dx$.\work{2.8cm}
  \item Find $\displaystyle\int\dfrac{x}{x+3}\,dx$.\work{2.8cm}
  \item Find $\displaystyle\int\dfrac{e^{2x}}{1+e^{4x}}\,dx$.\work{2.8cm}
  \item Find $\displaystyle\int\dfrac{\cos(\ln x)}{x}\,dx$ for $x>0$.\work{2.8cm}
\end{enumerate}

\begin{frbox}{Question 9 \textperiodcentered{} Free-response practice}
Evaluate each integral. Name the method you use.
\begin{enumerate}[label=(\alph*),leftmargin=*,itemsep=2pt]
  \item $\displaystyle\int_0^1\dfrac{x+2}{x^2+4x+1}\,dx$\work{2.6cm}
  \item $\displaystyle\int_0^1\dfrac{x}{x+1}\,dx$\work{2.6cm}
  \item $\displaystyle\int_0^{1/2}\dfrac{1}{\sqrt{1-x^2}}\,dx$\work{2.6cm}
\end{enumerate}
\end{frbox}

\newpage
\begin{solbox}
\begin{enumerate}[leftmargin=*,itemsep=4pt]
  \item Substitution $u=x^3+x$: $\dfrac{\left(x^3+x\right)^6}{6}+C$.
  \item Divide: $\displaystyle\int\left(1+x^{-2}\right)dx=x-\dfrac1x+C$.
  \item Arctangent with $u=2x$: $\dfrac12\arctan(2x)+C$.
  \item Substitution $u=e^{x}+2$: $\ln\!\left(e^{x}+2\right)+C$.
  \item $u=\cos x$: $\displaystyle\int-u^{-2}\,du=\dfrac1u=\sec x+C$.
  \item Division: $1-\dfrac{3}{x+3}$, giving $x-3\ln|x+3|+C$.
  \item $u=e^{2x}$, $du=2e^{2x}dx$: $\dfrac12\arctan\!\left(e^{2x}\right)+C$.
  \item $u=\ln x$: $\sin(\ln x)+C$.
  \item \begin{enumerate}[label=(\alph*),leftmargin=*,itemsep=3pt]
    \item Substitution $u=x^2+4x+1$, $du=2(x+2)\,dx$, limits 1 to 6: $\dfrac12\ln6$.
    \item Division: $\displaystyle\int_0^1\left(1-\dfrac{1}{x+1}\right)dx=1-\ln2$.
    \item Inverse sine: $\Big[\arcsin x\Big]_0^{1/2}=\dfrac\pi6$.
  \end{enumerate}
\end{enumerate}
\end{solbox}

\end{document}
